\documentclass[10pt,a4paper]{amsart}
\usepackage[margin=19mm]{geometry}
\usepackage{amsmath,amssymb,mathtools}
\usepackage[scr=rsfs,cal=euler]{mathalfa}
\usepackage{xfrac}
\usepackage[unicode,hidelinks,bookmarks=false]{hyperref}
\newcommand{\ie}{i.e.\ }
\newcommand{\cf}{cf.\ }
\newcommand{\resp}{respectively\ }
\newcommand{\Subsection}{\subsection}
\providecommand{\nobreakdash}{\nobreak\hbox{-}}
\setlength{\emergencystretch}{2em}
\multlinegap0pt
\usepackage{amssymb}
\usepackage{mathtools}
\usepackage[shortlabels]{enumitem}
\newtheorem{lemma}{Lemma}
\newtheorem{proposition}{Proposition}
\newcommand{\Id}{\mathrm{Id}}
\newcommand{\astar}{a_*}
\newcommand{\cI}{\mathcal I}
\newcommand{\one}{\mathbf 1}
\title{Entropy and singularity of spectral and pivotal sets in planar percolation}
\author{Yitzchak Shmalo}
\date{Source: arXiv:2609.15388v1 (CC BY 4.0)}
\begin{document}
\maketitle

\noindent\emph{Source and licence.} Entropy and singularity of spectral and pivotal sets in planar percolation. Version arXiv:2609.15388v1 (CC BY 4.0), distributed by arXiv under the Creative Commons Attribution 4.0 International licence, http://creativecommons.org/licenses/by/4.0/, as recorded in the arXiv licence metadata for this exact version and shown on the abstract page https://arxiv.org/abs/2609.15388v1. Full licence notice: \texttt{licenses/arxiv-2609.15388-CC-BY-4.0.txt}.\par\smallskip

\noindent\textit{Editorial scope.} Counted unit: the article's proposition prop:coercivity (source0.tex lines 872--878). The article's complete original proof of it is delivered with it (source0.tex lines 879--924, one \texttt{proof} environment of 46 lines). No mathematics is written, repaired or re-derived: the statement and the proof are the article's own lines, in the article's own notation, and no retained line is substituted or re-flowed.  Retained uncounted as the source-local context the proof needs: lines 163--165; lines 863--865.  Nothing was omitted from any retained span, so the package carries no omission record.  No retained line contains a citation, so the wrapper carries no bibliography and no bibliography entry.  No retained line contains a figure environment, an image-inclusion command or drawing code, so no image is delivered and the package declares no asset dependency.  Editorial formatting only, in addition: an AMS \texttt{amsart} preamble that restates the article's own declarations and macros used by the retained lines, namely the environments \texttt{lemma}, \texttt{proposition} on separate counters, together with the standard packages those lines need.  The retained environments print in this document as Lemma 1, Proposition 1 in order of appearance, whereas the article prints the same items with the numbering of its own full document.  The complete native source of the article as distributed in the arXiv e-print for this exact version is bundled unchanged as \texttt{sources/210412/source0.tex}.
\begin{equation}\label{eq:astar}
 \astar=\frac{8-\sqrt{53}}{16}>\frac1{32},
\end{equation}

\begin{lemma}\label{lem:rank}
If $q\in\operatorname{Ran}Q$ vanishes on $L$, then $q=0$.
\end{lemma}

\begin{proposition}\label{prop:coercivity}
Every function $g$ on six signs which vanishes on $L$ satisfies
\begin{equation}\label{eq:coercivity}
 \|(\Id-Q)g\|_2^2\ge\astar\|g\|_2^2.
\end{equation}
The constant in \eqref{eq:astar} is optimal on this function space.
\end{proposition}

\begin{proof}
Lemma~\ref{lem:rank} already implies the inequality with some positive constant by compactness of the unit sphere of functions supported on $H$. We compute the best constant.

Let $M$ be the $32\times18$ matrix
\[
 M_{z,A}=\chi_A(z),\qquad z\in L,\quad A\in\cI,
\]
and put $G=M^{\mathsf T}M$. If $q=\sum_{A\in\cI}u_A\chi_A$, then
\begin{equation}\label{eq:gramnorm}
 \|q\|_2^2=\|u\|^2,\qquad
 \|\one_Lq\|_2^2=\frac1{64}u^{\mathsf T}Gu.
\end{equation}
We record the small-block calculation to specify the constant without numerical approximation.

The symmetry $z\mapsto-z$ separates even and odd degrees. Order the even block by the constant and the six distance-two pairs, followed by the three opposite pairs. Then
\[
 G_{\rm even}=32I_{10}+
 \begin{pmatrix}0&K\\K^{\mathsf T}&0\end{pmatrix},
 \qquad K^{\mathsf T}K=512I_3+112J_3.
\]
Here $J_3$ is the all-ones matrix. The first row of $K$ is $(-4,-4,-4)$; the other rows are two copies of each permutation of $(-4,12,12)$. The eigenvalues of this block are
\[
 32\pm4\sqrt{53},\quad
 32\pm16\sqrt2\ \text{(each twice)},\quad
 32\ \text{(four times)}.
\]

In the odd block, the singleton matrix is circulant with first row $(32,12,0,-4,0,12)$; its eigenvalues are $52,48,16,12,16,48$. The alternating triples have block
\[
 \begin{pmatrix}32&-4\\-4&32\end{pmatrix}.
\]
A singleton couples with coefficient zero to its own parity triple and with coefficient $-4$ to the other. Only the constant and alternating singleton modes couple to these triples, giving the two blocks
\[
 \begin{pmatrix}52&-4\sqrt3\\-4\sqrt3&28\end{pmatrix},
 \qquad
 \begin{pmatrix}12&4\sqrt3\\4\sqrt3&36\end{pmatrix}.
\]
Thus the odd eigenvalues are $40\pm8\sqrt3$, $24\pm8\sqrt3$, and $48,16$ each twice. The least eigenvalue of $G$ is $32-4\sqrt{53}$.

By \eqref{eq:gramnorm}, $\|\one_Lq\|_2^2\ge\astar\|q\|_2^2$ on $\operatorname{Ran}Q$. If $g=\one_Hg$, then
\[
 |\langle g,q\rangle|=|\langle g,\one_Hq\rangle|
 \le\sqrt{1-\astar}\,\|g\|_2\|q\|_2.
\]
Take $q=Qg$ and use Pythagoras to obtain \eqref{eq:coercivity}. For sharpness, choose a least-eigenvalue vector $q$ in \eqref{eq:gramnorm} and put $g=\one_Hq$. Then $Qg=(1-\astar)q$, and equality follows.
\end{proof}

\end{document}
